\section{受禅之后：把北方守住，把江南接进来}
\label{sec:sui-reunification}

开皇元年二月，北周静帝把帝位让给杨坚，国号改为隋。受禅的仪式把一场已在禁军、中枢和关陇政治网络中发生的权力转移，写成了有次序的继承；它却没有自动消除北方的军事压力，也没有让长江以南立刻成为新朝的土地。新朝首先得到的是关中及其军政资源，随后才要回答一个更困难的问题：怎样把这批资源调度到边塞、法庭和江南州县。%
\footnote{受禅的月份和基本过程见《隋书·高祖纪上》《北史·隋本纪上》及《资治通鉴》开皇元年二月条；三书可互校年序，却都不能直接说明朝中每一位行动者的动机。}

这里的材料本身要求读者暂缓赞颂或谴责。《隋书》由唐初朝廷组织修成，后来的唐人需要解释自己何以继承隋，又何以不必为隋的覆亡负责；于是，隋文帝常被安排为能恢复秩序的开端，隋炀帝则容易成为结局的承担者。这种叙述并非没有保存事实，而是把事实编排进后继王朝的治乱图式。把《北史》与《资治通鉴》的年序一同放在桌上，可以较稳地确认受禅、南征和镇压的先后；至于朝廷内外何人真心拥戴、何人只是服从既成力量，现存文字远不如仪文整齐。

开皇二年，沙钵略可汗联结北周旧部南下，提醒隋廷：改朝换代首先会在边界上接受检验。隋的回应不是把北方问题当作一役即可毕结的战争，而是在守御之外经营分化和交涉，设法不让突厥汗国的压力与南方战事同时压到关中。汉文史书把这一过程写成新朝能够驾驭“归附”与“离叛”的故事；但汗国内部的联结、首领的选择及牧地上的计算，大多只能透过隋廷的报告看见。因而可以说，北方威胁暂未阻断隋的计划，却不能把它写成一个已经被征服、边界已经安静的世界。%
\footnote{《隋书·高祖纪上》《隋书·突厥传》与《北史·突厥传》共同构成这一段的编年骨架；它们主要保留中原朝廷的观察，未留下可与之对读的突厥方面连续叙事。}

边防尚在调整，朝廷也在给新国家寻找共同的办事语言。开皇三年修定律令，减省死刑条目并重整刑名。法律在这里不是装饰性的“文明”标志：北周旧律、关中与关东不同的治理经验，需要有一套可供官府援引的规范。开皇律后来成为唐律的重要前史，说明这次整理留下了长久的制度资源；但《隋书·刑法志》和《通典》所保存的，是颁行与追述的文字，不能替每一所州县的审判作证。条文到达何处，如何由地方官、书吏和被告理解，又在哪些地方被变通，仍须等待个案文书，不能由后来的法典谱系倒推。%
\footnote{《隋书·刑法志》《北史·刑法志》与《通典·刑法》记载开皇律令的修定；它们适于说明规范的形成，不足以量化地方执行。}

到开皇八年，隋廷认为可以南下。战争从来不是一条“北方强、南方弱”的简单公式：多路水陆军能够越过长江，靠的是北方整合、边防暂获缓冲以及陈朝军政财政困境同时出现。十月发动的灭陈战争，次年正月以隋军进入建康、陈后主被俘告终。建康城门的易手可以标记南北分立的终点，却不能把陈朝宫廷故事当作江南社会的全部反应。围绕被俘者和宫廷逸事的文字，最容易被亡国叙事磨成道德寓言；《陈书·后主纪》与隋方本纪、编年史的并读，至少使读者看见这是不同政权留下的、彼此并不完全同声的记录。%
\footnote{开战见《隋书·高祖纪下》《隋书·陈后主纪》及《资治通鉴》开皇八年十月条；建康陷落见《隋书·高祖纪下》《陈书·后主纪》及《资治通鉴》开皇九年正月条。兵力、路线和宫廷细节多出自官修战功或亡国叙事，不能按精确数字和戏剧场面复原。}

真正的接收从胜利之后才开始。隋要清查户口，改置州县，接管陈朝军民，并让新的赋役与官吏进入原陈诸州；这些动作把“统一”落实为名册、仓廪、差役和地方命令。开皇十年江南多地起事，杨素等率军镇压，正说明建康失守并不等于地方已经接受新秩序。把叛乱只说成遗民的怀旧，或只说成新政的苛暴，都比材料走得更远：我们知道中央在接收中遭遇强烈反应，却很难从军功叙事中逐户辨认参与者的处境、诉求和死亡数。%
\footnote{《隋书·高祖纪下》《隋书·杨素传》与《资治通鉴》开皇十年条确认起事和镇压；参与规模、动机与伤亡主要经由中央军功叙事保存，无法据此作精确统计。}

城址的分期、墓志所留下的姓名和日期，能够为建康及江南地方的转换补上少数独立锚点，也能提醒我们官修史书并非唯一入口；它们却不能把城市遗存自动变成全体居民的态度，把精英墓志变成普通家庭的账本。对隋初而言，考古材料最有力的作用，是约束过于整齐的叙事，而不是替我们补造一幅人人同样服从或人人同样反抗的图景。

因此，隋的再统一既不是一纸禅让的自然结果，也不只是一次渡江的胜负。它是防守北边、重整法律、调动水陆军和在江南反复建立州县控制的连锁过程。581年至590年的事件把帝国重新接合起来，同时也暴露出接合所需的征发、登记与武力。后来唐人以隋为鉴时，看见的往往是成败分明的前史；沿着这些较短、较难看的步骤看去，所见则是一个新国家怎样在不完整的证据中，一面取得服从，一面制造新的摩擦。
